\documentclass[a4paper,11pt]{book}
%\documentclass[a4paper,twoside,11pt,titlepage]{book}
\usepackage{listings}
\usepackage[utf8]{inputenc}
% \usepackage[style=list, number=none]{glossary} %
%\usepackage{titlesec}
%\usepackage{pailatino}
\usepackage[spanish]{babel}
\usepackage{amsmath,amssymb,amsthm,mathtools}
\usepackage{mdframed}
\usepackage{lipsum}
\usepackage{dcolumn}
\usepackage{titlesec}

\newcolumntype{.}{D{.}{\esperiod}{-1}}
\makeatletter
\makeatother

%%%%%%%%%% LO QUE YO HE INTRODUCIDO %%%%%%%%%%%%
 % Recuadrar teorema/lema y asignarlo a un capítulo
\newmdtheoremenv{teoremaBox}{Teorema}[chapter]
\newmdtheoremenv{lemaBox}{Lema}[chapter]
\newmdtheoremenv{proposicionBox}{Proposición}[chapter]

\newtheorem{observacion}{Observación}
 % Usar recuadro negro el terminar demostración
\renewcommand\qedsymbol{$ \blacksquare $}
%\renewcommand\proof{\textit{Demostración:}  \qedsymbol}

 % EVITAR QUE PONGA "CHAPTER *" AL INICIO DEL CAPÍTULO
\titleformat{\chapter}[display]{\normalfont\bfseries}{}{0pt}{\Huge}

\def\spanishoperators{adj traza vect dom dist sop vol sgn  Hess Jac rango grado diag img longitud Maximizar Minimizar Optimizar sec cotan cosec}
\newcommand{\refPar}[1]{(#1)}
\newcommand{\w}{\displaystyle}
\newcommand{\norm}[1]{\left\Vert#1\right\Vert}
\newcommand{\abs}[1]{\left\vert#1\right\vert}
\newcommand{\pre}[1]{\left\langle#1\right\rangle}
\newcommand{\set}[1]{\left\{#1\right\}}
\newcommand{\RR}{\mathbb{R}}
\newcommand{\NN}{\mathbb{N}}
\newcommand{\ZZ}{\mathbb{Z}}
\newcommand\restr[2]{{% ejemplo: \restr{f}{A}
		\left.\kern-\nulldelimiterspace #1 \vphantom{\big|} \right|_{#2} 
}}

\newcommand{\vecN}[2][N]{ $\{{#2}_{1},...,{#2}_{#1} \}$ }
\newcommand{\vecSpace}{ $V$ }

%%%%%%%%%%%%%%%%%%%%%%%%%%%%%%%%%%%%%%%%%%%%%%%%%%%%%%%%%%%%%%

%\usepackage[chapter]{algorithm}
\RequirePackage{verbatim}
%\RequirePackage[Glenn]{fncychap}
\usepackage{fancyhdr}
\usepackage{graphicx}
\usepackage{afterpage}

\usepackage{longtable}

\usepackage[pdfborder={000}]{hyperref} %referencia

% ********************************************************************
% Re-usable information
% ********************************************************************
\newcommand{\myTitle}{Título del proyecto\xspace}
\newcommand{\myDegree}{Grado en ...\xspace}
\newcommand{\myName}{Nombre Apllido1 Apellido2 (alumno)\xspace}
\newcommand{\myProf}{Nombre Apllido1 Apellido2 (tutor1)\xspace}
\newcommand{\myOtherProf}{Nombre Apllido1 Apellido2 (tutor2)\xspace}
%\newcommand{\mySupervisor}{Put name here\xspace}
\newcommand{\myFaculty}{Escuela Técnica Superior de Ingenierías Informática y de
Telecomunicación\xspace}
\newcommand{\myFacultyShort}{E.T.S. de Ingenierías Informática y de
Telecomunicación\xspace}
\newcommand{\myDepartment}{Departamento de ...\xspace}
\newcommand{\myUni}{\protect{Universidad de Granada}\xspace}
\newcommand{\myLocation}{Granada\xspace}
\newcommand{\myTime}{\today\xspace}
\newcommand{\myVersion}{Version 0.1\xspace}


\hypersetup{
pdfauthor = {\myName (email (en) ugr (punto) es)},
pdftitle = {\myTitle},
pdfsubject = {},
pdfkeywords = {palabra_clave1, palabra_clave2, palabra_clave3, ...},
pdfcreator = {LaTeX con el paquete ....},
pdfproducer = {pdflatex}
}

%\hyphenation{}


%\usepackage{doxygen/doxygen}
%\usepackage{pdfpages}
\usepackage{url}
\usepackage{colortbl,longtable}
\usepackage[stable]{footmisc}
%\usepackage{index}

%\makeindex
%\usepackage[style=long, cols=2,border=plain,toc=true,number=none]{glossary}
% \makeglossary

% Definición de comandos que me son tiles:
%\renewcommand{\indexname}{Índice alfabético}
%\renewcommand{\glossaryname}{Glosario}

\pagestyle{fancy}
\fancyhf{}
\fancyhead[LO]{\leftmark}
\fancyhead[RE]{\rightmark}
\fancyhead[RO,LE]{\textbf{\thepage}}
\renewcommand{\chaptermark}[1]{\markboth{\textbf{#1}}{}}
\renewcommand{\sectionmark}[1]{\markright{\textbf{\thesection. #1}}}

\setlength{\headheight}{1.5\headheight}

\newcommand{\HRule}{\rule{\linewidth}{0.5mm}}


%Definimos los tipos teorema, ejemplo y definición podremos usar estos tipos
%simplemente poniendo \begin{teorema} \end{teorema} ...
\newtheorem{teorema}{Teorema}[chapter]
\newtheorem{ejemplo}{Ejemplo}[chapter]
\newtheorem{definicion}{Definición}[chapter]

\definecolor{gray97}{gray}{.97}
\definecolor{gray75}{gray}{.75}
\definecolor{gray45}{gray}{.45}
\definecolor{gray30}{gray}{.94}

\lstset{ frame=Ltb,
     framerule=0.5pt,
     aboveskip=0.5cm,
     framextopmargin=3pt,
     framexbottommargin=3pt,
     framexleftmargin=0.1cm,
     framesep=0pt,
     rulesep=.4pt,
     backgroundcolor=\color{gray97},
     rulesepcolor=\color{black},
     %
     stringstyle=\ttfamily,
     showstringspaces = false,
     basicstyle=\scriptsize\ttfamily,
     commentstyle=\color{gray45},
     keywordstyle=\bfseries,
     %
     numbers=left,
     numbersep=6pt,
     numberstyle=\tiny,
     numberfirstline = false,
     breaklines=true,
   }
 
% minimizar fragmentado de listados
\lstnewenvironment{listing}[1][]
   {\lstset{#1}\pagebreak[0]}{\pagebreak[0]}

\lstdefinestyle{CodigoC}
   {
	basicstyle=\scriptsize,
	frame=single,
	language=C,
	numbers=left
   }
\lstdefinestyle{CodigoC++}
   {
	basicstyle=\small,
	frame=single,
	backgroundcolor=\color{gray30},
	language=C++,
	numbers=left
   }

 
\lstdefinestyle{Consola}
   {basicstyle=\scriptsize\bf\ttfamily,
    backgroundcolor=\color{gray30},
    frame=single,
    numbers=none
   }


\newcommand{\bigrule}{\titlerule[0.5mm]}


%Para conseguir que en las páginas en blanco no ponga cabecerass
\makeatletter
\def\clearpage{%
  \ifvmode
    \ifnum \@dbltopnum =\m@ne
      \ifdim \pagetotal <\topskip
        \hbox{}
      \fi
    \fi
  \fi
  \newpage
  \thispagestyle{empty}
  \write\m@ne{}
  \vbox{}
  \penalty -\@Mi
}
\makeatother

\usepackage{pdfpages}
\begin{document}
%\frontmatter
%\tableofcontents
%\listoffigures
%\listoftables
%
%\mainmatter
%\setlength{\parskip}{5pt}

%\begin{titlepage}
 
 
\newlength{\centeroffset}
\setlength{\centeroffset}{-0.5\oddsidemargin}
\addtolength{\centeroffset}{0.5\evensidemargin}
\thispagestyle{empty}

\noindent\hspace*{\centeroffset}\begin{minipage}{\textwidth}

\centering
\includegraphics[width=0.9\textwidth]{imagenes/logo_ugr.jpg}\\[1.4cm]

\textsc{ \Large TRABAJO FIN DE GRADO\\[0.2cm]}
\textsc{ DOBLE GRADO EN INGENIERÍA EN INGENIERÍA INFORMÁTICA Y MATEMÁTICAS}\\[1cm]
% Upper part of the page
% 
% Title
{\Huge\bfseries Titulo del Proyecto\\
}
\noindent\rule[-1ex]{\textwidth}{3pt}\\[3.5ex]
{\large\bfseries Subtitulo del Proyecto}
\end{minipage}

\vspace{2.5cm}
\noindent\hspace*{\centeroffset}\begin{minipage}{\textwidth}
\centering

\textbf{Autor}\\ {Pedro Manuel Flores Crespo (alumno)}\\[2.5ex]
\textbf{Directores}\\
{Manuel Ruíz Galán (tutor1)\\
Juan Carlos Torres Cantero (tutor2)}\\[2cm]
\includegraphics[width=0.3\textwidth]{imagenes/etsiit_logo.png}\\[0.1cm]
\textsc{Escuela Técnica Superior de Ingenierías Informática y de Telecomunicación}\\
\textsc{---}\\
Granada, mes de 201
\end{minipage}
%\addtolength{\textwidth}{\centeroffset}
%\vspace{\stretch{2}}
\end{titlepage}




\chapter{Teorema de Hanh-Banach y de Mazur-Orlicz}
	\thispagestyle{empty}
	\paragraph{}El objetivo principal de esta sección es demostrar tanto una versión no tan conocida del clásico teorema de Hanh-Banach como el teorema de Mazur-Orlicz. Para ello, iremos utilizando una serie de lemas previos que nos facilitarán el proceso. Así, una vez vistos estos resultados, habremos construido la base para el lema de Simons.
	
	\paragraph{} En primer lugar vamos a recordar la definición de funcional sublineal sobre un espacio vectorial \vecSpace . Notar que todos los espacios vectoriales que vamos a usar son reales. Del mismo modo, los espacios y conjuntos que usaremos asumiremos que son no vacíos.
	
	\begin{definicion}
		Sea \vecSpace un espacio vectorial. Decimos que el $P:\vecSpace \rightarrow \RR$ es sublineal si cumple las siguientes condiciones:
		\begin{itemize}
			\item $ P $ es subaditiva: $x_1, x_2 \in \vecSpace \Longrightarrow P(x_1 + x_2) \leq P(x_1) + P(x_2) $.
			\item $ P $ es positivamente homogénea: $x_1 \in \vecSpace $ y $ \lambda > 0 \Longrightarrow P(\lambda x) = \lambda P(x) $.
		\end{itemize}
	\end{definicion}

	Como consecuencia, podemos afirmar que $ P(0) = 0 $. En efecto:
	\[
	P(0) = P(2\times0) = 2\times P(0) \Longrightarrow P(0) = 0.
	\]
	\paragraph{}Por ejemplo, toda norma o seminorma sobre \vecSpace es un funcional sublineal. Así, dados $ a,b \in \RR $ con $ a <b $ tenemos que $ P:H^{1}_{0} (a,b) \longrightarrow R$ definida como $ P(x) = \norm{x'}_{L^2(a,b)} $ es sublineal.También, si $ \vecSpace = \RR $ y definimos la parte positiva $ P(x) = (x)^{+} = \max \{0,x\}, \forall x \in \RR $ obtenemos un funcional sublineal sobre $\RR$.
	
	\paragraph{} El lema que exponemos a continuación, y que generalizaremos posteriormente en el lema \ref{lema2}, nos servirá para demostrar el teorema de Hanh-Banach. 
	
	\begin{lemaBox}\label{lema1}
		Sea\vecSpace un espacio vectorial y $P:\vecSpace \rightarrow \RR$ un funcional sublineal. Fijamos un elemento $ y \in V $. Para todo $ x \in V $ tomamos  
		\begin{center}
			$ P_y(x) := \inf_{\lambda > 0} \left[P(x+\lambda y) - \lambda P(y)\right] $
		\end{center}
		
		Entonces, $ P_y:V \rightarrow \RR$, $ P_y $es sublineal, $ P_y \leq P $ y $ P_y (-y) \leq  -P(y)$.
	\end{lemaBox} 
	\begin{proof}
		Fijamos $ y \in V $. Sea $ x \in V $ y $ \lambda > 0$. Como P es sublineal tenemos: 
		\begin{center}
			$ \lambda P(y) = P(\lambda y) =P(\lambda y +x-x) \leq P(x+\lambda y)+ P(-x)$.
		\end{center}
		Por lo tanto, se obtiene que $ P(x+\lambda y) - \lambda P(y) \geq -P(-x) $.  Tomando el ínfimo sobre $ \lambda >0 $ llegamos a $ P_{y}(x)\geq -P(-x) > -\infty$. Por consiguiente, $ P_y:V \rightarrow \RR$. \\
		
		Probaremos ahora que $ P_y $ es sublineal. Empezamos viendo la subaditividad. Tomamos $ x_1, x_2 \in V $ y sean $ \lambda_1, \lambda_2 > 0$ arbitrarios. Entonces: 
		\begin{equation*}
		\begin{split}
		\left[ P(x_1 + \lambda_1 y) - \lambda_1 P(y) \right] &+ \left[ P(x_2 + \lambda_2 y) - \lambda_2 P(y) \right] \\
		& \geq \left[ P(x_1 + x_2 + (\lambda_1+\lambda_2)y) \right] - (\lambda_1+\lambda_2) P(y) \\
		&\geq P_y (x_1 + x_2 ).
		\end{split}
		\end{equation*}
		Tomando ínfimo sobre $ \lambda_1 $ y $ \lambda_2 $, $  P_y (x_1)  + P_y (x_2 ) \geq P_y (x_1 + x_2 ) $. Así, $ P_y $ es subaditiva. Para comprobar que es positivamente homogénea tomamos $ x \in V $ y $ \mu > 0 $. Entonces:
		\begin{equation*}
		\begin{split}
		P_y (\mu x) &= \inf_{\lambda > 0} \left[P(\mu x+\lambda y) - \lambda P(y)\right] = \mu \inf_{\lambda > 0} \left[P(x+ (\lambda / \mu) y) - (\lambda / \mu) P(y)\right] \\
		&= \mu \inf_{\upsilon > 0} \left[P(x+ \upsilon y) - \upsilon  P(y)\right] = \mu P_y (x).
		\end{split}
		\end{equation*}	
		Obtenemos que $ P_y $ es positivamente homogénea y como consecuencia sublineal. \\
		
		Para demostrar que $ P_y \leq P $, sea $ x \in V $ y tomando $ \lambda = 1 $: 
		\[ P_y(x) \leq P(x+y) - P(y) \leq P(x)+ P(y) - P(y) = P(x) \Longrightarrow P_y \leq P.  \]
		Como $ x \in \vecSpace $ es arbitrario, entonces $ P_y \leq P $. Finalmente, actuando de manera similar al caso anterior:
		\[P_y(-y) \leq P(-y+y) - P(y) = -P(y).  \]
		
	\end{proof}
	\paragraph{} Ahora procedemos a probar el teorema de Hanh-Banch para funcionales sublineales, el cual es uno de los resultados más importantes del análisis funcional.	
	\begin{teoremaBox}[Hanh-Banach]\label{H-B}
		Sea V un espacio vectorial y $P:V \rightarrow \RR$ un funcional sublineal. Entonces existe un funcional lineal L en V tal que $ L \leq P $.
	\end{teoremaBox}
	\begin{proof}
		Sea $ SUB = \lbrace Q:\vecSpace \longrightarrow \RR : Q \leq P \rbrace$, es decir, el conjunto de funcionales sublineales sobre \vecSpace que son menores o iguales que $ P $. Nuestro propósito es emplear el lema de Zorn. Para ello, dados $ T_1 ,T_2 \in \mathcal{Q} $ consideramos la relación de orden usual, es decir:
		\begin{center}
			$ T_1 \leq T_2 \Longleftrightarrow T_1 (x) \leq T_2 (x) \quad \forall x \in \vecSpace $.
		\end{center}
		Primero probaremos que todo subconjunto $ \mathcal{T} $ totalmente ordenado de $SUB $ tiene una cota inferior en $ SUB $. Definimos $ Q(x):=\inf \{ T(x): T \in \mathcal{T} \} $ y queremos ver que $ Q(x) \in \RR $. Si $ x \in \vecSpace $ y $ T \in \mathcal{T} $, como T es subaditiva obtenemos la siguiente relación:
		\[ 0 = T(0) = T(x-x) \leq T(x) + T(-x) \Longrightarrow T(x) \geq -T(-x) \quad (1)\]  
		Por otro lado:
		\[ T \in \mathcal{Q} \Longrightarrow T(x) \leq P(x) \Longrightarrow -T(x) \geq -P(x) \quad (2)\] 
		Usando (1), (2) y tomando ínfimo sobre $  T $  llegamos a $ Q(x) \geq -P(x) \geq - \infty $. Por lo tanto $ Q:V \rightarrow \RR$. \\
		
		Ahora probaremos que $ Q $ es subaditiva. Para ello, tomamos $ x_1, x_2\in \vecSpace $. Sean $ T_1 , T_2 \in \mathcal{T} $ arbitrarios. Si $ T_1 \geq T_2 $ (el caso de $ T_2 \geq T_1 $ es análogo.):
		
		\begin{center}
			$ T_1 (x_1)+  T_2 (x_2) \geq T_2(x_1)+  T_2 (x_2) \geq T_2(x_1 +x_2) \geq Q(x_1 + x_2).$
		\end{center}
		Concluimos (en ambos casos) que $ T_1 (x_1)+  T_2 (x_2) \geq Q(x_1 + x_2)$. Tomando ínfimo en $ T_1 $ y $ T_2 $ obtenemos que $ Q (x_1)+  Q(x_2) \geq Q(x_1 + x_2)$. Así, $ Q $ es sublineal. Que sea positivamente homogénea es consecuencia de que $ T $ también lo es. Dado $ \mu > 0 $:
		\begin{equation*}
		\begin{split}
		Q(\mu x) &=\inf \{ T(\mu x): T \in \mathcal{T} \} \\ 
		& = \inf \{ \mu T( x): T \in \mathcal{T} \} \\ 
		&= \mu\inf \{ T( x): T \in \mathcal{T} \} \\ 
		&= \mu Q(x). 
		\end{split}
		\end{equation*}
		
		De este modo, Q es sublineal y como es claro que $ Q \leq P \Longrightarrow Q \in SUB$. Así, es directo que $ Q $ es el elemento minimal de $ \mathcal{T} $ en $SUB$.\\
		
		El lema de Zorn nos proporciona entonces un elemento minimal de $ SUB $ que llamaremos $ L $. Vamos a comprobar que $ L $ es lineal y, por tanto, es el funcional buscado. Tomamos ahora $ y \in \vecSpace $. Con la notación del lema anterior, $ L_y : \vecSpace \longrightarrow \RR $ es sublineal, $ L_y \leq L $ (como consecuencia $ L_y \in SUB $) y $ L_y (-y) \leq L(-y) $. De hecho, como $ L $ es minimal en $ SUB $, $ L_y = L $ y por ello $ L (-y) \leq L(-y) $. Por otro lado, como L es subaditiva, $ L(-y) \geq -L(y) $. Combinando ambas desigualdades, $ L(-y) = -L(y) $. Tomamos $ x \in \vecSpace $ y $ \lambda < 0 $, usando la igualdad anterior llegamos a:
		\[ \qquad \quad
		L(\lambda x) = L (-(-\lambda)x) = -L(-\lambda x) = -(-\lambda)L(x) = \lambda L(x) \label{1}
		\] 
		obteniendo que $ L $ es homogénea. Si $ x_1, x_2 \in \vecSpace $, la subaditividad de $ L $ nos da $ L(-x_1-x_2) \leq L(-x_1) + L(-x_2) $. Usando la homogeneidad de $ L $:
		\begin{equation*}
		\begin{split} \qquad
		L(x_1+x_2) &= L(-(-x_1-x_2)) = -L(-x_1-x_2) \\ 
		& \geq -L(-x_1)-L(-x_2) = L(x_1) + L (x_2) \geq L(x_1+x_2). 
		\end{split}
		\end{equation*}
		Por ello, $	L(x_1+x_2) = L(x_1) + L (x_2) $ y por la arbitrariedad de $ x_1, x_2 \in \vecSpace $ concluimos que $ L $ es lineal.
		
	\end{proof}

	\paragraph{} El siguiente resultado importante que demostraremos será el teorema de Mazur-Orlicz. Primero, veamos un lema previo.
	
	\begin{lemaBox}\label{lema2}
		Sea\vecSpace un espacio vectorial y $P:\vecSpace \rightarrow \RR$ un funcional sublineal. Sea $ D $ un subconjunto convexo de \vecSpace y sea $ \beta := \inf_D P \in \RR $. Para todo $ x \in V $ tomamos  
		\begin{center}
			$ Q(x) := \inf_{d \in D, \lambda > 0} \left[P(x+\lambda d) - \lambda \beta\right] $
		\end{center}
		
		Entonces, $ Q:V \rightarrow \RR$, $ Q $es sublineal, $ Q \leq P $ y $ \forall d \in D, -Q(-d) \geq \beta$.
	\end{lemaBox} 
	\begin{proof}
		Si $ x \in \vecSpace,\text{ }d \in D $ y $ \lambda > 0 $, entonces
		\begin{center}
			$ P(x+ \lambda d) - \lambda \beta \geq -P(-x) + \lambda P(d)-\lambda\beta \geq -P(-x) \geq -\infty$.
		\end{center}
		La primera igualdad se deduce de la linealidad de P ya que:
		\begin{equation*}
		\begin{split}
		\lambda P(d) &= P(\lambda d) \\ 
		&=P(\lambda d +x-x) \\ 
		&\leq P(x+\lambda d)+ P(-x)
		\end{split}
		\end{equation*}
		por lo que $ \lambda P(d)-P(-x) \leq P(x+\lambda d) $. La segunda se debe a que
		\[\beta = \inf_D P \Longrightarrow \lambda P(d) \geq \lambda\beta \Longrightarrow\lambda P(d) - \lambda\beta \geq 0. \]
		Tomando el ínfimo sobre $ d \in D  $ y $ \lambda > 0 $ llegamos a \[ Q(x)\geq -P(-x) > -\infty\] por lo que $ Q:V \rightarrow \RR$. Es relativamente fácil probar que $ Q $ es positivamente homogénea por lo que para ver que es sublineal solo queda comprobar la subaditividad. Para ello, tomamos $ x_1, x_2 \in V $. Sean $ d_1, d_2 \in D $ y $ \lambda_1, \lambda_2 > 0$ arbitrarios. Para simplificar la notación llamamos $ x := x_1 + x_2 $, $ \lambda := \lambda_1 + \lambda_2 $ y $ d:= (\lambda_1/\lambda)d_1 + (\lambda_2/\lambda)d_2 $. Notar que $ d \in D $ al ser este convexo. Entonces: 
		\begin{equation*}
		\begin{split}
		\left[ P(x_1 + \lambda_1 d_1) - \lambda_1 \beta \right] + \left[ P(x_2 + \lambda_2 d_2) - \lambda_2 \beta \right] &\geq P(x + \lambda_1 d_1 + \lambda_2 d_2) - \lambda \beta\\
		& = P(x +\lambda d) - \lambda \beta\\ 
		& \geq Q(x) = Q(x_1 + x_2)
		\end{split}.
		\end{equation*}
		Tomando ínfimo sobre $ d_1, d_2, \lambda_1 $ y $ \lambda_2 $, $  Q(x_1) + Q (x_2 ) \geq Q (x_1 + x_2 ) $. Así, $ Q $ es subaditiva y como consecuencia sublineal. Fijamos $ d \in D $. Sea $ x \in \vecSpace $ arbitrario. Entonces, $ \forall \lambda > 0 $, $ Q(x) \leq P(x) + \lambda \left[P(d) - \beta \right]$. Tomando $ \lambda \longrightarrow 0 $, $ Q(x) \leq P(x)$ y como consecuencia $ Q \leq P $. Finalmente, sea $ d \in D $ arbitrario y tomando $ \lambda = 1 $:
		\[Q(-d) \leq Q(-d+d) - \beta = -\beta \Longrightarrow  -Q(-d) \geq \beta.  \]
		
	\end{proof}

	\paragraph{}Visto este lema, estamos preparados para ver el resultado que nos interesa:
	
	\begin{teoremaBox}[Mazur-Orlicz]\label{M-O}
		Sea\vecSpace un espacio vectorial y $P:\vecSpace \rightarrow \RR$ un funcional sublineal.  Sea $ D $ un subconjunto no vacío y convexo de \vecSpace. Entonces existe un funcional $ L: \vecSpace \longrightarrow \RR $ lineal tal que $ L \leq P $ e $ \inf_D L = \inf_D P $.
	\end{teoremaBox}
	\begin{proof}
		Sea $ \beta := \inf_D P $. En el caso de que $ \beta = -\infty $ por el teorema de Hanh-Banach tenemos que $ \exists L $ sobre \vecSpace tal que es lineal y $ L \leq P$. Así:
		\begin{center}
			$ L \leq P \Longrightarrow inf_D L \leq \inf_D P = -\infty \Longrightarrow inf_D L = \inf_D P.$ 
		\end{center}
		Supongamos entonces que $ \beta \in \RR $. Definimos el funcional auxiliar $ Q $ como en el lema \ref{lema2}. Del teorema de Hanh-Banach obtenemos que existe un funcional lineal $ L $ sobre \vecSpace tal que $ L \leq Q$ (como $ Q \leq P $ tenemos que $ L \leq P $). Sea $ d \in D $, entonces:
		\[
		L(d) = -L(-d) \geq -Q(-d) \geq \beta.
		\]
		Tomando ínfimo sobre $ d \in D $:
		\[
		\inf_D L \geq \beta = \inf_D P.
		\]
		Por otro lado, como $ L \geq P $:
		\[
		\inf_D L \leq\inf_D P.
		\]
		Juntando ambas desigualdades obtenemos $ \inf_D L =\inf_D P $.
	\end{proof}



\chapter{Igualdad}
\newcommand{\normSpace}{E}
	
Ahora, queremos obtener una igualdad que nos será de utilidad posteriormente para el Teorema de Separación. Recordemos dados $ \normSpace_1, \normSpace_2 $ dos espacios normados y $ L:\normSpace_1 \longrightarrow \normSpace_2 $ un operador lineal, entonces L es continuo si, y solo si, verifica la siguiente condición:
\[
\exists \alpha > 0 : \norm{L(x)} \leq \alpha \norm{x} \quad \forall x \in \normSpace_1,
\]
conocido como el espacio dual (topológico) de $ E $.
Consideramos el espacio vectorial dado por:
\[
\normSpace^* =  \lbrace L:\normSpace \longrightarrow \RR: L \text{ es lineal y continuo} \rbrace.
\]
Para todo $ L \in \normSpace^* $ definimos su norma como:
\[
\norm{L} = \min \lbrace \alpha > 0:  \norm{L(x)} \leq \alpha \norm{x} \quad \forall x \in \normSpace\rbrace.
\]
De este modo, podemos escribir:
\begin{equation*}\label{desigNorma}
	\norm{L(x)} \leq \norm{L} \norm{x}
\end{equation*}
siendo dicha desigualdad óptima. También podemos expresar su norma como el mínimo mayorante de un conjunto mayorado, es decir, el supremo:
\[
\norm{L} = \sup \lbrace \norm{L(x)}/\norm{x} : \forall x \in \normSpace \setminus \{0\}\rbrace.
\]
Para $  x \in \normSpace_1 \setminus \{0\} $ tenemos que $ \norm{L(x)}/\norm{x} = \norm{L(x/\norm{x})}$ y es claro que $ \{ x/\norm{x} : x \in \normSpace_1 \setminus \{0\}\} $ es la esfera unidad de $ \normSpace $ que notamos como $ S_\normSpace $. Si en vez de la esfera consideramos la bola unidad, $ B_\normSpace $ el supremo no varía. Efectivamente, si $ x \in B_\normSpace  $ se tiene que $ x = \norm{x}u $ con $ u \in  S_\normSpace$, y por ello $ \norm{L(x)} = \norm{x}\norm{L(u)} \leq \norm{L(u)} $ ya que $ \norm{x} \leq 1 $. De este modo, también tenemos que:
\begin{equation*}\label{normSup}
	\norm{L} =\sup_{x \in B_\normSpace} \norm{L(x)}.
\end{equation*}

En este momento, estamos en disposición de enunciar y demostrar la igualdad que deseamos:
\begin{proposicionBox}
	Dado un espacio normado $ \normSpace $ y $ x \in \normSpace $, entonces se cumple que:
	\begin{equation}\label{iguNorm}
	\sup_{x^* \in B_{\normSpace ^ *}} x^*(x) = \norm{x}.
	\end{equation}
\end{proposicionBox}
\begin{proof}
Sea $ x_0 \in E$ y sea $ D := \{x_0\} $ Consideramos el funcional 
\begin{equation*}
\begin{split}
P:\normSpace \longrightarrow &\RR \\
x \longmapsto &\norm{x}.
\end{split}
\end{equation*} 
Es claro que $ P $ es sublineal y que $ D $ es convexo. Podemos aplicar el Teorema de Mazur-Orlicz, teorema \ref{M-O}, y obtenemos que existe un funcional $ L:\normSpace \longrightarrow \RR $ lineal tal que $ L \leq P $ e $ \inf_D L = \inf_D P $. Como $ L \leq P $, entonces 
\begin{equation*}
 	\norm{L(x)} \leq \norm{P(x)} = \norm{x} \quad \forall x \in \normSpace.
\end{equation*}
Concluimos que $ L $ es continua. Además, como tenemos que $ L \in \normSpace ^* $, llamamos $ L = x^* $ y llegamos a que $ \norm{L} = \norm{x^*}  \leq 1$, es decri, $ x^* \in  B_{\normSpace ^ *} $.  Por su parte, como  $ \inf_D L = \inf_D x^* = \inf_D P $ y $ D = \{x_0\} $ entonces, $ x^*(x_0) = \norm{x_0} $. De este modo, llegamos a que existe $ x^* \in B_{\normSpace^*} $ tal que $ x^*(x_0) = \norm{x_0} $. Si tomamos cualquier elemento $ y^* \in B_{\normSpace^*}  $, entonces $ \norm{y^*(x_0)} \leq \norm{y^*} \norm{x_0} \leq \norm{x_0} $ y podemos asegurar que:
\[
\sup_{x^* \in B_{\normSpace ^ *}} x^*(x_0) = \norm{x_0}.
\] 
Como $ x_0 \in \normSpace $ es arbitrario, la desigualdad enunciada queda probada.
\end{proof}

\chapter{Lema de Simons y Teorema de la Alternativa de Gordan}
	\thispagestyle{empty}
	\newcommand{\ttt}{\textbf{\emph{t}}}
	\newcommand{\sss}{\textbf{\emph{s}}}
	
	\paragraph{}Este capítulo se centra en lema de Simons, el cual será de gran ayuda para demostrar el Teorema de la Alternativa de Gordan.

	\paragraph{}Antes de comenzar, necesitamos hacer la siguiente definición.
	
	\begin{definicion}
		Dado $ N \in \NN $ con $ N > 1 $ llamamos símplex unitario de $ \RR^N $ al conjunto convexo y compacto de $ \RR^N $ dado por:
		\begin{equation*}
			\Delta_N := \left\lbrace \ttt \in \RR^N: \sum_{i=1}^{N}{t_i} = 1 \hspace{0.5em} , \hspace{0.5em} t_1,...,t_N \geq 0 \right\rbrace .
		\end{equation*}
	\end{definicion}

	Por ejemplo, si $ N = 2 $, tenemos que $ \Delta_2 = \mathrm{co} \{(1,0),(0,1)\}$, es decir es la envolvente convexa de los vectores de la base usual de $ \RR^2 $. De hecho en general, $ \Delta_N = co\{e_1,...e_N\} $. \\
	
	Antes de continuar veamos que efectivamente $ \Delta_N $ es convexo y compacto:
	
	\begin{itemize}
		\item  Convexo: tenemos que comprobar que dados $ \ttt, \sss \in \Delta_N  \Longrightarrow \lambda\ttt + (1-\lambda)\sss \in \Delta_N, \lambda \in \lbrack 0,1 \rbrack$. En efecto, las coordenadas de $ \lambda\ttt + (1-\lambda)\sss \in \RR^N $ verifican:
		\begin{itemize}
			\item [i) ] $ \lambda t_i + (1-\lambda)s_i \geq 0 $ para todo $ i = 1,..., N $ ya que $ t_i, s_i \geq 0 $.
			\item [ii) ] $ \sum_{i=1}^{N} (\lambda t_i + (1-\lambda)s_i) =   \lambda \sum_{i=1}^{N} t_i + (1-\lambda)\sum_{i=1}^{N} s_i = \lambda + (1 - \lambda) = 1$ ya que $ \sum_{i=1}^{N} t_i = \sum_{i=1}^{N} s_i = 1  $.
		\end{itemize}
	
		Por lo tanto, $ \lambda\ttt + (1-\lambda)\sss \in \Delta_N \text{ para todo }\lambda \in \lbrack 0,1 \rbrack $.
		
		\item Compacto: al encontrarnos en $ \RR^N $ y aplicando el conocido Teorema de Heine-Borel basta ver que $ \Delta_N $ es cerrado y acotado. Claramente es acotado por lo que nos centraremos en la de cerrado. Sea $ \lbrace \ttt_n \rbrace_{ n \in \NN} $ una sucesión de $ \Delta_N $ y sea $ \ttt_0 \in \RR^N $ tal que $ \lbrace \ttt_n \rbrace_{ n \in \NN} \longrightarrow \ttt_0 $. Tenemos que comprobar que $ \ttt_0 \in \Delta_N $. 
		\begin{itemize}
		\item[i) ] Como todas las coordenadas de cada $ \ttt_n $ para $ n \in \NN $ son no negativas podemos asegurar que las de $ \ttt_0 $ también lo son.
		\item[ii) ] La función $ f: \RR^N \longrightarrow \RR $ definida como $ f(\ttt) =  \sum_{i=1}^{N} t_i $ es continua. Claramente $ f(\ttt) = 1 $ para todo $ \ttt \in \Delta_N $ y por ello $ \lbrace f(\ttt_n) \rbrace_{ n \in \NN} \longrightarrow 1$. Por continuidad de $ f $ y unicidad de límite tenemos que $ f(\ttt_0) = 1 $ pero eso implica que la suma de sus componentes vale 1
		\end{itemize} 
	
		Así, hemos demostrado que $ \ttt_0 \in \Delta_N $ y por lo tanto $ \Delta_N $ es compacto. 
	\end{itemize}

	\paragraph{} Enunciamos este lema previo que usaremos posteriormente el la demostración del de Simons. 
	\begin{lemaBox}\label{lema2.1}
		Sea $ N \in \NN $, con $ N \geq 1 $ y $ L: \RR^N \longrightarrow \RR $ definida por $ S(x):=\max\{x_1,...,x_N\}$. Entonces, $ S $ es sublineal. Además, si $ L:\RR^N \longrightarrow \RR $ es un funcional lineal tal que $ L \leq S $ entonces $ L $ es de la forma $ L(x) = t_1 x_1 + ... + t_N x_N $ con $ (t_1,...,t_N) \in \Delta_N$. De hecho, el recíproco también es cierto, es decir, si $ L =  (t_1,...,t_N) \in \Delta_N $ entonces $ L \leq S $.
		
		%%%%%% NOTA %%%%%%%%
		%% Las aplicaciones lineales coninciden con el dual del espacio. Es decir, podemos identificar la aplicación con sus coeficientes obeniendo de ese modo un punto del espacio.
		
	\end{lemaBox}
	\begin{proof}
		Claramente, $ S $ es positivamente homogénea. También es subaditiva ya que dados $ x,y \in \vecSpace $ : 
		\begin{equation*}
		\begin{split}
		S(x+y) &= \max \{x_1 + y_1, ..., x_N +y_N \}\\ 
		&\leq \max\{x_1, ..., x_N\} + \max \{y_1, ...,y_N\} = S(x) + S(y)
		\end{split}
		\end{equation*}
		
		Por ello, $ S $ es sublineal. Para terminar veamos que 
		\begin{equation*}
			\left\lbrace L:\RR^N \longrightarrow \RR : L \text{ lineal y }L\leq S \right\rbrace = \Delta_N
		\end{equation*}
		
		\begin{itemize}
			\item[$ \supseteq $ )] Sea $ \mathbf{t} \in \Delta_N$ definimos $ L:\RR^N \longrightarrow \RR $ como $ L(x):= \langle \mathbf{t},x \rangle $. Es evidente que $ L $ es lineal al serlo el producto escalar. Dado $ x \in \RR $:
			\begin{equation*}
				L(x) = \sum_{i=1}^{N}{t_i x_i} \leq \sum_{i=1}^{N}{t_i S(x)} = S(x)\sum_{i=1}^{N}{t_i} = S(x)
			\end{equation*}
			donde la primera desigualdad se debe a que $ x_i \leq S(x)$ para todo $ x_i$ con $ i=1,...,N $ y a que $ t_i \geq 0 $ ya que $ \mathbf{t} \in \Delta_N$. Esto también justifica la que última igualdad ya que $ \sum_{i=1}^{N}{t_i} = 1 $.
			
			\item[$ \subseteq $ )] Sea $ L =  (t_1,...,t_N) \in \RR^N $ lineal tal que $ \forall x \in \RR^N $ cumple $ \sum_{i=1}^{N}{t_i x_i} \leq \max\{x_1,...,x_N\}$. Así, si tomamos $ e_i \in \RR^N $ donde $ e_i $ representa el i-ésimo elemento de la base usual de $ \RR^N $ con $ i=1,...,N $. Entonces:   		
			\begin{equation*}
				L(-e_i) = -t_i \leq 0 \Longrightarrow t_i \geq 0 \quad \forall i=1,...,N
			\end{equation*}
			Si ahora llamamos $ e = \sum_{i=1}^{N}{e_i} $. Obtenemos:
			\begin{equation*}
				\begin{rcases*}
				L(e) = \sum_{i=1}^{N}{t_i} \leq \max \{1, ...,1\} = 1 \\
				L(-e) = -\sum_{i=1}^{N}{t_i} \leq \max \{-1, ...,-1\} = -1
				\end{rcases*} \Longrightarrow \sum_{i=1}^{N}{t_i} = 1
			\end{equation*}
			Concluimos entonces que $ L =  (t_1,...,t_N) \in \Delta_N $.
		\end{itemize}
		
	\end{proof}

	\paragraph{} Enunciamos ahora lema de Simons\cite{Simons2008} en que destacamos la ausencia de hipótesis topológicas, lo que será importante posteriormente.
	
	\begin{lemaBox}[Lema de Simons]\label{Simons}
		Sea $ C $ un subconjunto no vacío y convexo de un espacio vectorial. Dadas $ f_1, ..., f_N $ ($ N \geq 1 $)  funciones sobre $ C $ reales y convexas, entones existen $ (t_1, ..., t_N) \in \Delta_N $ que cumplen
		\[
		\inf_C\left[ \max \{f_1, ..., f_N\}\right] = \inf_C \left[ t_1 f_1+ ...+t_N f_N \right]
		\] 
	\end{lemaBox}
	\begin{proof}
		Sea  \vecSpace = $ \RR^N $ con $ N \in \NN $. Definimos $S:\vecSpace \longrightarrow P $ como \[ S(x_1, ..., x_N) := \max \{x_1, ..., x_N\} \]. Por el lema \ref{lema2.1} $ S $ es sublineal. Tomamos el subconjunto:
		\[ 
		D = \{ (x_1, ..., x_N)\in V: \exists c \in C \quad \text{ tal que } \quad \forall i = 1,...N,\quad f_i(c) \leq x_i \}
		\]
		
		$ D $ es un subconjunto convexo de \vecSpace. Sean $ x, y \in D $, por ello, existen $ c_x, c_y \in C $ tales que $ f_i (c_x) \leq x_i  $ y $ f_i (c_y) \leq y_i \quad \forall i=1,...,N $. Dado $ \lambda \in [0,1] $, llamamos $ c := (1-\lambda)c_x + \lambda c_y $ que pertenece a $ C $ por ser este convexo. Veamos que $ c $ es el elemento necesario de $ C $ para que cualquier combinación convexa de $ x $ e $ y $ esté en $ D $. Así, para todo $ i =1,...,N  $:	
		\[
		f_i(c) = f_i((1-\lambda)c_x + \lambda c_y) \leq (1-\lambda)f(c_x) + \lambda f(c_y) \leq (1-\lambda)x_i + \lambda y_i 
		\]
		donde la primera desigualdad se debe a que las $ f_i $ son convexas y la segunda a que $ x,y \in D $. Por ello, $ (1-\lambda)x_i + \lambda y_i \in D , \quad \forall \lambda \in [0,1] $ por lo que $ D $ es convexo. Aplicando el Teorema de Mazur-Orlizc, existe $ L $ funcinal lineal sobre \vecSpace tal que $ L \leq S $ e $ \inf_D L = \inf_D S $. \\
		
		Nuevamente, por el lema \ref{lema2.1} tenemos que $ L = (t_1,...,t_N) \in \Delta_N$. \\
		
		Finalmente:
		\[
		\inf_D L = \inf_{c\in C} \left[ t_1 f(c) + ..+ t_N f(c) \right] = \inf_{C} \left[ t_1 f + ..+ t_N f \right]
		\]
		y
		\[
		\inf_D S = \inf_{c\in C} \left[ \max \{f_1(c), ..., f_N(c)\} \right] = \inf_C\left[ \max \{f_1, ..., f_N\}\right] 
		\]
		por lo que 
		\[ \inf_{C} \left[ t_1 f + ..+ t_N f \right] = \inf_C\left[ \max \{f_1, ..., f_N\}\right]  \] 
	\end{proof}

	\paragraph{}Enunciamos ahora el Teorema de la alternativa de Gordan en su versión convexa.
	
	\begin{teoremaBox}[Teorema de la Alternativa de Gordan-versión convexa]\label{Gordan}
		Sea $ C $ un subconjunto convexo de un espacio vectorial y sean $ f_1,...,f_N : C \longrightarrow \RR $ funciones convexas ($ N \geq 1 $). Entonces una, y solo una, de la siguientes afirmaciones se cumple:
		\begin{itemize}
			\item[i)] $ \exists \mathbf{t} \in \Delta_N $ tal que $ 0 \leq \inf_{C}  \sum_{i=1}^{N}{t_i f_i}$.
			\item[ii)] $ \exists c \in C $ que cumple $ \max \lbrace f_1(c), ..., f_N(c) \rbrace < 0 $.
		\end{itemize}
	\end{teoremaBox}
	\begin{proof}
		Si aplicamos el lema de Simons, lema \ref{Simons} a las funciones $ f_1,...,f_N $ obtenemos:
		
		\begin{equation*}
			\exists t \in \Delta_N \text{ : } \inf_{ C}\left[ \max \lbrace f_1, ..., f_N \rbrace \right] = \inf_{C} \sum_{i=1}^{N}t_i f_i
		\end{equation*}
		
		Supongamos en primer lugar que $ \alpha := \inf_{C}\left[ \max \lbrace f_1, ..., f_N \rbrace \right] \in \RR $. Planteamos la siguiente alternativa, cuyos casos son excluyentes:
		\begin{itemize}
			\item[a)] $ 0 \leq \alpha $: implica $ i) $ ya que $ \alpha = \inf_{C} \sum_{i=1}^{N}t_i f_i $ 
			\item[b)] $ \alpha < 0 $: este caso, por su parte, implica $ ii) $ ya que:
			
			\begin{equation*}
				\alpha = \inf_{C}\left[ \max \lbrace f_1, ..., f_N \rbrace \right] < 0 \Longleftrightarrow \exists c \in C  : \max \lbrace f_1(c), ..., f_N(c) \rbrace < 0 
			\end{equation*}  
		\end{itemize}
	
		Para finalizar, veamos cuando $ \alpha =-\infty $. En este caso, estamos en la misma situación que en b) por lo que solo se puede dar $ ii) $.
	\end{proof}

	\paragraph{} Destacamos las siguientes observaciones:
	
	\begin{observacion}
		Esta versión convexa del teorema implica la versión clásica del mismo.
	\end{observacion}

	Dados $ \{x_1,...x_N\}$ con $ x_i \in \RR^M , (M \geq 1 )$ $i=1,...,N$, la versión clásica del teorema nos aporta las siguientes alternativas:

	\begin{itemize}
		\item[i*)] $ \exists \mathbf{t} \in \Delta_N $ tal que $ 0 = \sum_{i=1}^{N}{t_i x_i}$.
		\item[ii*)] $ \exists y \in \RR^M $ tal que cumple $ \max_{i=1,...,N} \langle y, x_i \rangle < 0 $.
	\end{itemize}

	Para ello, basta aplicar la versión convexa del teorema a $ C := \RR^M $ y a las funciones $ f_1,...,f_N : C \longrightarrow \RR $ definidas por $ f_i(c):=\langle c,x_i \rangle , \forall i=1,...,N  $. Notar que las funciones $ f_1,...,f_N $ son lineales por la izquierda y como consecuencia son convexas. En este caso, la alternativa $ ii) $ implica $ ii*) $ ya que:
	
	\begin{equation*}
		\exists c \in C = \RR^M \text{ : } \max_{i=1,...,N} {\langle c,x_i \rangle}  =  \max_{i=1,...,N}f_i (c) < 0 
	\end{equation*}
	
	Por su parte, la alternativa $ i) $ nos da:
	\begin{equation*}
		\exists \mathbf{t} \in \Delta_N \text{ : } 0 \leq \inf_{c \in \RR^M}  \sum_{i=1}^{N}{t_i f_i(c) } = \inf_{c \in \RR^M} \sum_{i=1}^{N}{t_i\langle c,x_i \rangle} = \inf_{c \in \RR^M} \langle c, \sum_{i=1}^{N}{t_i x_i} \rangle 
	\end{equation*}
	
	Hemos obtenido por ello que  $0  \leq \inf_{c \in \RR^M} \langle c, \sum_{i=1}^{N}{t_i x_i} \rangle  $ lo que nos lleva a $ 0 \leq \langle c, \sum_{i=1}^{N}{t_i x_i} \rangle  $ para todo $ c \in \RR^M $. Usando la linealidad por la izquierda del producto escalar:
	\[
	0 \leq \langle -c, \sum_{i=1}^{N}{t_i x_i} \rangle \Longleftrightarrow 	0 \leq -\langle c, \sum_{i=1}^{N}{t_i x_i} \rangle 
	\Longleftrightarrow  \langle c, \sum_{i=1}^{N}{t_i x_i}\rangle \leq 0, \quad \forall c \in \RR^M
	\]
	
	Juntando ambas desigualdades obtenemos que $ 0 =  \langle c, \sum_{i=1}^{N}{t_i x_i}\rangle, \quad \forall c \in \RR^M $. Como la igualdad anterior se cumple para todo elemento de $ \RR^M $ entonces podemos deducir que $ \sum_{i=1}^{N}{t_i x_i} = 0 $	ya que $ \sum_{i=1}^{N}{t_i x_i} \in (\RR^M)^{\perp} = \{0\} $. Así pues, tenemos que
	\begin{equation*}
		\text{Se cumple }i) \Longleftrightarrow \exists t \in \Delta_N \text{ : }  0 = \sum_{i=1}^{N}{t_i x_i}
	\end{equation*} 	
	
	Es claro entonces que obtenemos $ i*) $.
	
	\begin{observacion}
		El lema de Simons(lema \ref{Simons}) y el Teorema de la Alternativa de Gordan (teorema \ref{Gordan}) son equivalentes.
	\end{observacion}

	\paragraph{} Ya hemos visto que el Lema de Simons implica el Teorema de la Alternativa de Gordan. Veamos que el recíproco también es cierto.  \\
	
	Llamamos $ \alpha := \inf_{C}\left[ \max \lbrace f_1, ..., f_N \rbrace \right] $. Si $ \alpha = -\infty $. Por el lema \ref{lema2.1} sabemos que $ \forall \ttt \in \Delta_N $ se cumple que $ \sum_{i=1}^{N} t_i f_i(c) \leq \max \lbrace f_1 (c), ... , f_N (c) \rbrace$ para todo $ c \in C$. Tomando ínfimos en C:
	\[
	\inf_C\left[ \sum_{i=1}^{N} t_i f_i \right] \leq \inf_C \left[ \max \lbrace f_1, ..., f_N \rbrace \right] = -\infty \Longrightarrow \inf_C\left[ \sum_{i=1}^{N} t_i f_i \right] = -\infty 
	\]
	
	y por ello $ \forall \ttt \in \Delta_N $ (en particular para uno) se cumple que
	
	\[
	\inf_C\left[ \sum_{i=1}^{N} t_i f_i \right] = \inf_C \left[ \max \lbrace f_1, ..., f_N \rbrace \right] \]

	Supongamos ahora que $ \alpha \in \RR $. Sean las funciones $ q_1, ..., g_N: C \longrightarrow \RR $ definidas como $ g_i = f_i - \alpha $ con $ i=1,...,N$. Veamos que las funciones $ g_1, ..., g_N $ son convexas como consecuencia de que $ f_1, ..., f_N $ lo son. Sean $ c_1,c_2 \in C $ y $ \lambda \in \left[0,1\right] $:
	\begin{equation*}
	\begin{split}
	g_i(\lambda c_1 + (1-\lambda) c_2) &= f_i(\lambda c_1 + (1-\lambda) c_2) - \alpha \\
	&\leq \lambda f_i(c_1) + (1-\lambda)f_i(c_2) - \alpha \\
	&= \lambda f_i(c_1) + (1-\lambda)f_i(c_2) - \lambda \alpha + (1-\lambda)\alpha \\
	&= \lambda( f_i(c_1) - \alpha ) + (1-\lambda) (f_i (c_2) - \alpha) \\
	&= \lambda g_i(c_1) + (1-\lambda) g_i (c_2)
	\end{split}
	\end{equation*}
	
	Obtenemos que $ g_i $ es convexa para todo $ i = 1, ..., N $. Si usamos el Teorema de la Alternativa de Gordan obtenemos que solo se pueden dar una y solo de las siguientes posibilidades:
	
	\begin{itemize}
		\item[i)] $ \exists \mathbf{t} \in \Delta_N $ tal que $ 0 \leq \inf_{C}  \sum_{i=1}^{N}{t_i g_i}$.
		\item[ii)] $ \exists c \in C $ que cumple $ \max \lbrace g_1 (c), ..., g_N(c) \rbrace < 0 $.
	\end{itemize}
	
	Razonemos que no se puede dar $ ii) $. Si fuese así, tendríamos que $ \exists c \in C $ tal que $ \max \lbrace g_1(c), ..., g_N(c) \rbrace  =  \max \lbrace f_1(c) - \alpha, ..., f_N(c) - \alpha \rbrace < 0 $. En particular, existiría un índice $ j \in {1,...,N} $ que cumpliría $ f_j(c) - \alpha < 0 \Longrightarrow f_j(c) < \alpha = inf_{C}\left[ \max \lbrace f_1, ..., f_N \rbrace \right] $. Esto es imposible por la propia definición de ínfimo. Por ello, afirmamos que $ \exists \mathbf{t} \in \Delta_N $ tal que $ 0 \leq \inf_{C}  \sum_{i=1}^{N}{t_i g_i}$. Desarrollando el sumatorio:
	\begin{equation*}
	\begin{split}
	0 &< \inf_{C}  \sum_{i=1}^{N}{t_i g_i} = \inf_{C}  \sum_{i=1}^{N}{t_i(f_i - \alpha)} = \inf_{C} \left[ \sum_{i=1}^{N}{t_i(f_i)} - \sum_{i=1}^{N}{t_i\alpha} \right]\\
	&= \inf_{C} \left[ \sum_{i=1}^{N}{t_i(f_i)} -\alpha \sum_{i=1}^{N}{t_i} \right] = \inf_{C} \left[ \sum_{i=1}^{N}{t_i(f_i)} - \alpha \right] = \inf_{C} \left[ \sum_{i=1}^{N}{t_i(f_i)}\right] - \alpha
	\end{split}
	\end{equation*}
	
	Por lo tanto:
	\[
	0 < \inf_{C} \left[ \sum_{i=1}^{N}{t_i(f_i)}\right] - \alpha \Longleftrightarrow \inf_{C}\left[ \max \lbrace f_1, ..., f_N \rbrace \right] = \alpha  < \inf_{C} \left[ \sum_{i=1}^{N}{t_i(f_i)}\right]
	\]
	
	El lema \ref{lema2.1} nos aporta la otra desigualdad y llegamos nuevamente a que $ \exists \ttt \in \Delta_N $ que cumple:
		\[
	\inf_C\left[ \sum_{i=1}^{N} t_i f_i \right] = \inf_C \left[ \max \lbrace f_1, ..., f_N \rbrace \right] \]
	

\chapter{Teoremas de Fritz John y de Karush-Kuhn-Tucker}
		\newcommand{\barx}{\bar{x} }
		\newcommand{\dd}{\textbf{\emph{d}}}
		
	En primer lugar, empezamos recordando la definición de derivada direccional. 
	\begin{definicion}
			Sea $ D \subset \RR^M $ con $ M \in \NN $ y sea la función $ g: D \longrightarrow  \RR$, definimos la derivada direccional de g en $ x \in D $ en la dirección del vector $ d\in \RR^M $ como
			\[
			g'(x;d) = \lim_{t\rightarrow0}\frac{g(x+td) - g(x)}{t}
			\]
			siempre y cuando el límite exista. Diremos que g es diferenciable en el sentido de Gâteaux en x si $ g'(x;\cdot):\RR^M \longrightarrow \RR $ es lineal y en ese caso escribimos $ \nabla g(x) = g'(x;\cdot) $, es decir, $ g'(x;d) = \langle \nabla g(x), d\rangle $ con $ d \in \RR^M $.
	\end{definicion}
	
	Dentro del contexto de este trabajo, cuando decimos que una función es diferenciable nos referimos a que lo es en el sentido de Gâteaux. A estas funciones también las llamaremos Gâteaux diferenciables. Destacamos que este concepto de diferenciabilidad es más débil que el de Fréchet, que es el más frecuente dentro de nuestros estudios. De hecho, si una función $ g $ es diferenciable en $ x_0 \in D$ en el sentido de Fréchet, y notamos su derivada como $ Dg(x_0) $ entonces $ g $ también es diferenciable en el sentido de Gâteaux en $ x_0 $ y además $ Dg(x_0) = \nabla g(x_0) $. El recíproco no es cierto tal y como mostramos en el siguiente ejemplo:\\
	
	EJEMPLO \\
	
	A continuación, demostremos cómo se calcula la derivada de la función máximo, lo que nos será útil en posteriores resultados.
	\begin{proposicionBox}\label{dirDeriv}
		Sean $ D \subset \RR^M$ ($ M \in \NN $) , $ \barx $ un punto del interior de $ D $ y sean $ g_1, ..., g_N : D \longrightarrow \RR $  funciones continuas y diferenciables en $ \barx $ donde $ N \in \NN $. Definimos $ g:D \longrightarrow \RR $ como $ g(x):=\max_{i=1,...,N}\{g_i(x)\} $ y el conjunto de índices $ K = \lbrace  i : g_i(\barx) =  g(\barx) \rbrace $. Entonces, para toda dirección $ d \in \RR^M $ la derivada direccional de $ g $ existe en todo $ \RR^M $ y viene dada por la siguiente expresión:
		\begin{equation}
			g'(\barx;d) = \max_{i \in K}\{ \langle \nabla g_i(\barx), d \rangle\}, \quad d\in\RR^M.
		\end{equation}
	\end{proposicionBox}
	\begin{proof}
		Podemos suponer sin pérdida de generalidad que el conjunto $ K = \{1, ..., N \} $ ya que aquellas $ g_i $ que no alcancen el máximo no afectarán al cálculo de la derivada de $ g $. Para cada $ i \in K $ tenemos la siguiente desigualdad:
		\begin{equation*}
			\liminf_{t\rightarrow0}\frac{g(\barx+td) - g(\barx)}{t} \geq \liminf_{t\rightarrow0}\frac{g_i(\barx+td) - g_i(\barx)}{t} = \langle \nabla g_i(\barx), d \rangle.
		\end{equation*}	
		La primera desigualdad se deduce de la definición de $ g $ ya que es el máximo de las $ g_i $ para $ i=1,...,N$ y la segunda igualdad de que todas las $ g_i $ son diferenciables en $ \barx $ y por tanto existe el límite de la definición de derivada direccional y coincide con el límite inferior. Por lo tanto:
		\[
		\liminf_{t\rightarrow0}\frac{g(\barx+td) - g(\barx)}{t} \geq \max_{i=1,...,N}\langle \nabla g_i(\barx), d \rangle.
		\]
		Por otro lado, afirmamos que :
		\begin{equation*}
			\limsup_{t\rightarrow0}\frac{g(\barx+td) - g(\barx)}{t} \leq \max_{i=1,...,N}\langle \nabla g_i(\barx), d \rangle.
		\end{equation*}
		De lo contrario, existirían una sucesión $ \{t_n\}\rightarrow 0 $ y $ \varepsilon > 0 $ que cumplirían:
		\[
		\frac{g(\barx+t_n d) - g(\barx)}{t_n} \geq \max_{i=1,...,N}\langle \nabla g_i(\barx), d \rangle + \varepsilon \quad \forall n \in \NN.
		\]
		Tomamos ahora una sucesión parcial $ \{t_{\sigma(n)}\}_{n \in \NN} $ con $ \sigma:\NN \longrightarrow \NN $ estrictamente creciente y $ j \in K $ un índice fijo tal que para todo $ k \in \{\sigma(n)\}_{n \in \NN} $ se cumple que $ g(\barx + t_k d) = g_j (\barx + t_k d)$. Tomando límite obtenemos que :
		\begin{equation*}
		\begin{split}
		\limsup_{t\rightarrow0}\frac{g(\barx+td) - g(\barx)}{t} &= 	\limsup_{t\rightarrow0}\frac{g_j(\barx+td) - g_j(\barx)}{t} \\
		&= \langle \nabla g_j(\barx), d \rangle \\ &\geq \max_{i=1,...,N}\langle \nabla g_i(\barx), d \rangle + \varepsilon,
		\end{split}
		\end{equation*}
		lo cual, es imposible. Finalmente, hemos obtenido que:
		\[
		\limsup_{t\rightarrow0}\frac{g(\barx+td) - g(\barx)}{t} \leq \max_{i=1,...,N}\langle \nabla g_i(\barx), d \rangle \leq 	\liminf_{t\rightarrow0}\frac{g(\barx+td) - g(\barx)}{t}.
		\]
		Como el límite inferior es siempre menor o igual que el superior concluimos que ambos coinciden y por lo tanto existe el límite y además:
		\[
		\lim_{t\rightarrow0}\frac{g(\barx+td) - g(\barx)}{t} = g'(\barx;d) = \max_{i=1,...,N}\langle \nabla g_i(\barx), d \rangle.
		\]
	\end{proof}

	Nuestro objetivo ahora es encontrar soluciones a problemas del siguiente tipo:
	
		\begin{equation}\label{probMin}
		\begin{cases}
		\inf_{x\in D} f(x)\\
		\begin{split}
		\text{s.a. } g_1(x) &\leq 0 \\
		&\vdots \\
		g_N(x) &\leq 0, \quad N \in \NN
		\end{split}
		
		\end{cases} 
		\end{equation}
		donde $ D \subset \RR^M$, $ f  $ es la función objetivo y las restricciones $ g_i $ con $ i =1,\dots, N $ funciones reales definidas en $ D $ y continuas. Si un punto satisface todas las restricciones diremos que es \textit{factible}  y como consecuencia llamamos \textit{región de factibilidad} al conjunto de todos los puntos factibles. Para un punto factible $ \barx $ definimos el \textit{conjunto activo} como $ I(\barx) = \{i: g_i(\barx) = 0\}$. Para este problema y asumiendo que $ \barx \in C $, llamamos \textit{vector de multiplicadores de Lagrange para $ \barx $} a $ \lambda \in (\RR^N)^+ $ si $ \barx $ es un punto crítico de:
		\[
		L(x;\lambda) = f(x) + \sum_{i=1}^{N} \lambda_i g_i(x),
		\]
		es decir, se cumple que (cuando $ f,g_1,\dots,g_N $ sean diferenciables):
		\[
		\nabla f(\barx) + \sum_{i=1}^{N} \lambda_i \nabla g_i(\barx) = 0
		\]
		y además $ \lambda_i = 0 $ si $ i \notin I(\barx) $.
		\begin{teoremaBox}[Teorema de Fritz John]\label{FritzJohn}
			Supongamos que el problema (\ref{probMin}) tiene un mínimo local en $ \barx \in D $. Si las funciones $ f, g_i $ con $ i \in I(\barx) $ son diferenciables en $ \barx $ entonces existen $ \lambda_0, \lambda_i \in \RR^+ $ para $ i \in I(\barx) $, no todas cero, que satisfacen:
			\[
			\lambda_0 \nabla f(\barx) + \sum_{i \in I(\barx)} \lambda_i \nabla g_i(\barx) = 0.
			\]
		\end{teoremaBox}
		\begin{proof}
			Consideramos la función 
			\[
			g(x) = \max \{ f(x) - f(\barx),\text{ } g_i(x) : i \in I(\barx)\}.
			\]
			Como $ \barx $ es un mínimo local del problema (\ref{probMin}) también lo es de $ g $. Esto se debe a que como $ f(x) \geq f(\barx) $ entonces $ f(x) - f(\barx) \geq 0$ para todo $ x $ en un entorno de $ \barx $. Por otro lado, $ g_i(x) \leq 0 $ para todo $ x \in D $ y como $ i \in I(\barx) $ entonces $ g_i(\barx)=0 $. De este modo $ g(x) \geq 0 \text{ } $ para todo $ x $ en un entorno de $ \barx $ y $ g(\barx) = 0 $ por lo que efectivamente alcanza un mínimo local en $ \barx $. Por la proposición \ref{dirDeriv} tenemos que para toda dirección $ d \in \vecSpace $ se cumple:
			\[
			g'(\barx;d) = \max \{ \langle \nabla f(\barx), d\rangle , \langle \nabla g_i(\barx),d \rangle : i \in I(\barx)\} \geq 0,
			\]
			ya que si $ g'(\barx;d) < 0 $, para todo $ t > 0 $ suficientemente pequeño tendríamos que $ g(\barx + td) < g(\barx) $ lo que contradice que $ g $ alcanza de mínimo local en $ \barx $. \\
			
			Por lo tanto, el sistema 
			\begin{equation*}
			\begin{cases}
			\langle \nabla f(\barx), d\rangle  < 0 \\
			\langle \nabla g_i(\barx),d \rangle < 0\text{ con } i \in I(\barx)
			\end{cases}
			\end{equation*}
			
			no tiene solución (para ninguna dirección) ya que al menos uno es no negativo. Si aplicamos el Teorema de la Alternativa de Gordan en su versión clásica, teorema \ref{GordanClasic}, vemos que solo se puede dar la alternativa i*) y en ese caso obtenemos que:
			\[
			 \exists \ttt = (t_0, ..., t_{M})\in \Delta_{M+1}  \text{ tal que }0 = t_0 \nabla f(\barx) + \sum_{i \in I(\barx)}  t_i \nabla g_i (\barx),
			 \]
			con $ M $ el cardinal del conjunto $ I(\barx) $. La demostración concluye llamando $ \lambda_0 = t_0 $ y $ \lambda_i = t_i $ con $ i \in I(\barx) $.
		\end{proof}
	
		\paragraph{}El teorema de Fritz John nos pueden aportan una gran desventaja y es que es posible que $ \lambda_0 = 0 $ por lo que la función objetivo es independiente a las restricciones. Por ello, necesitamos imponer algunas condiciones extra. En esta situación diremos que se cumple la \textit{condición de Mangasarian-Fromovitz} si existe una dirección $ d_0 \in \vecSpace $ que satisface que $ \langle \nabla g_i(\barx),d_0 \rangle < 0 $ para todo índice $ i \in I(\barx)$. Enunciamos ahora otro teorema que soluciona el problema que acabamos de comentar.
		
		\begin{teoremaBox}[Teorema de Karush-Kuhn-Tucker]
		Supongamos que el problema (\ref{probMin}) tiene un mínimo local en $ \barx \in D $. Si las funciones $ f, g_i $ con $ i \in I(\barx) $ son diferenciables en $ \barx $ y se cumple la condición de Mangasarian-Fromovitz entonces existe un vector de multiplicadores de Lagrange para $ \barx $.
		\end{teoremaBox} 
	
		\begin{proof}
		Del teorema \ref{FritzJohn} de las condiciones de Fritz John obtenemos que:
		\[
		\exists \ttt = (t_0, ..., t_{M})\in \Delta_{M+1}  \text{ tal que }0 = t_0 \nabla f(\barx) + \sum_{i \in I(\barx)}  t_i \nabla g_i (\barx)
		\]		
		con $ M $ el cardinal de $ I(\barx) $. Si multiplicamos escalarmente la igualdad por $ d_0 $ (dirección del espacio vectorial que nos aporta el requisito de  Mangasarian-Fromovitz) obtenemos:
		\[
		0 = t_0 \langle \nabla f(\barx), d_0 \rangle + \sum_{i \in I(\barx)}  t_i \langle \nabla g_i (\barx), d_0 \rangle.
		\]
		El requisito también nos da que $ t_0 \neq 0$. Razonemos por reducción al absurdo. Si se diese el caso tendríamos que
		\[
		0 = \sum_{i \in I(\barx)}  t_i \langle \nabla g_i (\barx), d_0 \rangle.
		\]
		Al tener $ \ttt \in \Delta_{M+1} $ se cumple que todas sus compnentes son positivas y $ \sum_{i \in I(\barx)}  t_i = 1 $ (estamos suponiendo que $ t_0 = 0 $ por lo que no influye en la suma de la definición de $ \Delta_{M+1} $) por lo que algún término es distinto de 0. Tenemos garantizado que $  \langle g_i (\barx), d_0 \rangle < 0 \quad \forall i \in I(\barx) $. Así tendríamos que:
			\[
		0 = \sum_{i \in I(\barx)}  t_i \langle \nabla g_i (\barx), d_0 \rangle < 0,
		\]
		lo cual es imposible. Por ello, concluimos que $ t_0 \neq 0 $. La demostración concluye tomando $ \lambda_i = t_i / t_0 $ para $ i \in I(\barx) $.
		\end{proof}
		

\chapter{Minimax}
		\newcommand{\barx}{\bar{x} }
		\newcommand{\dd}{\textbf{\emph{d}}}
	\paragraph{}
	
	\begin{definicion}
			Sea $ D $ un subconjunto de un espacio vectorial real y $ g: D \longrightarrow  \RR$ definimos la derivada direccional de g en la dirección del vector d del espacio vectorial como
			\[
			g'(x;d) = \lim_{t\rightarrow0}\frac{g(x+td) - g(x)}{t}
			\]
			siempre y cuando el límite exista. Diremos que g es diferenciable en el sentido de Gâteaux en x si $ g'(x;d) = \langle \nabla g(x), d\rangle $.
	\end{definicion}

	Dentro de este contesto, cuando decimos que una función es diferenciable nos referimos a que es diferenciable en el sentido de Gâteaux. Antes de continuar demostremos cómo se calcula la derivada de la función máximo, lo que nos será útil en posteriores resultados.
	
	\begin{proposicionBox}\label{dirDeriv}
		Sea $ D $ un subconjunto de un espacio vectorial real, $ \barx $ un punto del interior de $ D $ y $ g_1, ..., g_N : D \longrightarrow \RR $ funciones continuas y diferenciables en $ \barx $. Definimos $ g:D \longrightarrow \RR $ como $ g(x):=\max_{i=1,...,N}\{g_i(x)\} $ y el conjunto de índices $ K = \lbrace  i : g_i(\barx) =  g(\barx) \rbrace $. Entonces, para toda dirección del espacio vectorial, llamémosla d, la derivada direccional de $ g $ viene dada por la siguiente expresión:
		\begin{equation}
			g'(\barx;d) = \max_{i \in K}\{ \langle \nabla g_i(\barx), d \rangle\}
		\end{equation}
	\end{proposicionBox}
	\begin{proof}
		Podemos suponer sin pérdida de generalidad que $ K = \{1, ..., N \} $ ya que aquellas $ g_i $ que no alcancen el máximo no afectarán al cálculo de la derivada de $ g $. Para cada $ i \in K $ tenemos la siguiente desigualdad:
		\begin{equation*}
			\liminf_{t\rightarrow0}\frac{g(\barx+td) - g(\barx)}{t} \geq \liminf_{t\rightarrow0}\frac{g_i(\barx+td) - g_i(\barx)}{t} = \langle \nabla g_i(\barx), d \rangle
		\end{equation*}
		
		La primera desigualdad se deduce de la definición de $ g $ ya que es el máximo de las $ g_i $ para $ i=1,...,N$ y la segunda igualdad a que todas las $ g_i $ son diferenciables en $ \barx $ y por tanto existe el límite de la definición de derivada direccional y coincide con el límite inferior. Por lo tanto:
		\[
		\liminf_{t\rightarrow0}\frac{g(\barx+td) - g(\barx)}{t} \geq \max_{i=1,...,N}\langle \nabla g_i(\barx), d \rangle
		\]
		
		Por otro lado, afirmamos que :
		\begin{equation*}
			\limsup_{t\rightarrow0}\frac{g(\barx+td) - g(\barx)}{t} \leq \max_{i=1,...,N}\langle \nabla g_i(\barx), d \rangle
		\end{equation*}
		
		De lo contrario, existirían una sucesión $ \{t_n\}\rightarrow 0 $ y $ \varepsilon > 0 $ que cumplirían:
		\[
		\frac{g(\barx+t_n d) - g(\barx)}{t_n} \geq \max_{i=1,...,N}\langle \nabla g_i(\barx), d \rangle + \varepsilon \quad \forall n \in \NN
		\]
		
		Tomamos ahora una sucesión parcial $ \{t_{\sigma(n)}\}_{n \in \NN} $ con $ \sigma:\NN \longrightarrow \NN $ estrictamente creciente y $ j \in K $ un índice fijo tal que para todo $ k \in \{\sigma(n)\}_{n \in \NN} $ se cumple que $ g(\barx + t_k d) = g_j (\barx + t_k d)$. Tomando límite obtenemos que :
		\begin{equation*}
		\begin{split}
		\limsup_{t\rightarrow0}\frac{g(\barx+td) - g(\barx)}{t} &= 	\limsup_{t\rightarrow0}\frac{g_j(\barx+td) - g_j(\barx)}{t} \\
		&= \langle \nabla g_j(\barx), d \rangle \geq \max_{i=1,...,N}\langle \nabla g_i(\barx), d \rangle + \varepsilon \quad
		\end{split}
		\end{equation*}

		lo cual es imposible. Finalmente, hemos obtenido que:
		\[
		\limsup_{t\rightarrow0}\frac{g(\barx+td) - g(\barx)}{t} \leq \max_{i=1,...,N}\langle \nabla g_i(\barx), d \rangle \leq 	\liminf_{t\rightarrow0}\frac{g(\barx+td) - g(\barx)}{t}
		\]
		
		Como el límite inferior es siempre menor o igual que el superior concluimos que ambos coinciden y por lo tanto existe el límite y además:
		\[
		\lim_{t\rightarrow0}\frac{g(\barx+td) - g(\barx)}{t} = g'(\barx;d) = \max_{i=1,...,N}\langle \nabla g_i(\barx), d \rangle
		\]
	\end{proof}

	Nuestro objetivo ahora es encontrar soluciones a problemas del siguiente tipo:
	
		\begin{equation}\label{probMin}
		\begin{cases}
		\text{minimizar } f(x)\\
		\text{sujeto a } g_i(x) \leq 0 \quad \forall i =1,...,N\\
		\text{con }x \in C
		\end{cases} 
		\end{equation}
		
		donde $ C $ es un subconjunto del espacio vectorial, $ f  $ es la función objetivo y las restricciones $ g_i $ con $ i \in I $ son continuas en $ C $. Si un punto satisface todas las restricciones diremos que es \textit{factible}  y como consecuencia llamamos \textit{región de factibilidad} al conjunto de todos los puntos factibles. Para un punto factible $ \barx $ definimos el \textit{conjunto activo} como $ I(\barx) = \{i: g_i(\barx) = 0\}$. Para este problema y asumiendo que $ \barx \in C $ llamamos \textit{vector de multiplicadores de Lagrange para $ \barx $} a $ \lambda \in (\RR^N)^+ $ si $ \barx $ es un punto crítico de:
		\[
		L(x;\lambda) = f(x) + \sum_{i=1}^{N} \lambda_i g_i(x)
		\]
		
		\begin{teoremaBox}[Condiciones de Fritz John]\label{FritzJohn}
			Supongamos que el problema \ref{probMin} tiene un mínimo local en $ \barx \in C $. Si las funciones $ f, g_i $ con $ i \in I(\barx) $ son diferenciables en $ \barx $ entonces existen $ \lambda_0, \lambda_i \in \RR^+ $ para $ i \in I(\barx) $, no todas cero, que satisfacen:
			\[
			\lambda_0 \nabla f(\barx) + \sum_{i \in I(\barx)} \lambda_i \nabla g_i(\barx) = 0
			\]
		\end{teoremaBox}
		\begin{proof}
			Consideramos la función 
			\[
			g(x) = \max \{ f(x) - f(\barx),\text{ } g_i(x) : i \in I(\barx)\}
			\]
			
			Como $ \barx $ es un mínimo local del problema \ref{probMin} también lo es de $ g $. Esto se debe a que como $ f(x) \geq f(\barx) \Longrightarrow  f(x) - f(\barx) \geq 0$ para todo $ x $ en un entorno de $ \barx $. Por otro lado, $ g_i(x) \leq 0 $ si $ x \in C $ y si $ i \in I(\barx) $ entonces $ g_i(\barx)=0 $. De este modo $ g(x) \geq 0 \text{ } $ para todo $ x $ en un entorno de $ \barx $ y $ g(\barx) = 0 $ por lo que efectivamente alcanza un mínimo local en $ \barx $. Por la proposición \ref{dirDeriv} tenemos que para toda dirección d del espacio vectorial se cumple:
			\[
			g'(\barx;d) = \max \{ \langle \nabla f(\barx), d\rangle , \langle \nabla g_i(\barx),d \rangle : i \in I(\barx)\} \geq 0
			\]
			
			ya que si $ g'(\barx;d) < 0 $, para todo $ t > 0 $ suficientemente pequeño tendríamos que $ g(\barx + td) < g(\barx) $ lo que contradice la existencia de mínimo local en $ \barx $.
			
			Por lo tanto, el sistema 
			\begin{equation*}
			\begin{cases}
			\langle \nabla f(\barx), d\rangle  < 0 \\
			\langle \nabla g_i(\barx),d \rangle < 0\text{ con } i \in I(\barx)
			\end{cases}
			\end{equation*}
			
			no tiene solución (para ninguna dirección) ya que al menos uno es no negativo. Si aplicamos el Teorema de la Alternativa de Gordan en su versión clásica (referenciar ??) vemos que solo se puede dar la alternativa i*) y en ese caso obtenemos que:
			\[
			 \exists \ttt = (t_0, ..., t_{M})\in \Delta_{M+1}  \text{ tal que }0 = t_0 \nabla f(\barx) + \sum_{i \in I(\barx)}  t_i \nabla g_i (\barx)
			 \]
			 
			con $ M $ el cardinal del conjunto $ I(\barx) $. La demostración concluye llamando $ \lambda_0 = t_0 $ y $ \lambda_i = t_i $ con $ i \in I(\barx) $
		\end{proof}
	
		Las condiciones de Fritz John nos pueden aportan una gran desventaja y es que es posible que $ \lambda_0 = 0 $ por lo que la función objetivo es independiente a las restricciones. Por ello, necesitamos imponer algunas condiciones extra. En esta situación diremos que se cumple el \textit{requisito de Mangasarian-Fromovitz} si existe una dirección d del espacio vectorial que satisface que $ \langle \nabla g_i(\barx),d \rangle < 0 $ para todo índice $ i \in I(\barx)$. Enunciamos ahora otro teorema que soluciona el problema que acabamos de comentar.
		
		\begin{teoremaBox}[Condiciones de Karush-Kuhn-Tucker]
		Supongamos que el problema \ref{probMin} tiene un mínimo local en $ \barx \in C $. Si las funciones $ f, g_i $ con $ i \in I(\barx) $ son diferenciables en $ \barx $ y se cumple el requisito de Mangasarian-Fromovitz entonces existe un vector de multiplicadores de Lagrange para $ \barx $.
		\end{teoremaBox} 
	
		\begin{proof}
		El teorema \ref{FritzJohn} de las condiciones de Fritz John nos aporta que:
		\[
		\exists \ttt = (t_0, ..., t_{M})\in \Delta_{M+1}  \text{ tal que }0 = t_0 \nabla f(\barx) + \sum_{i \in I(\barx)}  t_i \nabla g_i (\barx)
		\]
		
		con $ M $ el cardinal de $ I(\barx) $. Si multiplicamos escalarmente la igualdad por $ d_0 $ obtenemos:
		\[
		0 = t_0 \langle \nabla f(\barx), d_0 \rangle + \sum_{i \in I(\barx)}  t_i \langle \nabla g_i (\barx), d_0 \rangle
		\]
		
		El requisito de Mangasarian-Fromovitz nos aporta que $ t_0 \neq 0$. Razonemos por reducción al absurdo. Si se diese el caso tendríamos que
		\[
		0 = \sum_{i \in I(\barx)}  t_i \langle \nabla g_i (\barx), d_0 \rangle
		\]
		
		Al tener $ \ttt \in \Delta_{M+1} $ se cumple que todas sus compnentes son positivas y $ \sum_{i \in I(\barx)}  t_i = 1 $ por lo que algún término es distinto de 0. El requisito de Mangasarian-Fromovitz garantiza que $  \langle g_i (\barx), d_0 \rangle < 0 \quad \forall i \in I(\barx) $. Así tendríamos que:
			\[
		0 = \sum_{i \in I(\barx)}  t_i \langle \nabla g_i (\barx), d_0 \rangle < 0
		\]
		lo cual es imposible. Por ello, concluimos que $ t_0 \neq 0 $.
		\end{proof}
		

%% Incluir en la bibliografía las referencias no citadas
% \nocite{Alfonso2010a}
\bibliographystyle{alpha-es}
\nocite{borwein2010convex}
\bibliography{bibliografia/bibliografia}

%\bibliographystyle{apalike-es}
%\bibliography{bibliografia/bibliografia}
%
%\input{capitulos/03_Planificacion}
%
%\input{capitulos/04_Analisis}
%
%\input{capitulos/05_Diseno}
%
%\input{capitulos/06_Implementacion}
%
%\input{capitulos/07_Pruebas}
%
%\input{capitulos/08_Conclusiones}
%
%%\chapter{Conclusiones y Trabajos Futuros}
%
%
%%\nocite{*}
%\bibliography{bibliografia/bibliografia}\addcontentsline{toc}{chapter}{Bibliografía}
%\bibliographystyle{miunsrturl}
%
%\appendix
%\input{apendices/manual_usuario/manual_usuario}
%%\input{apendices/paper/paper}
%\input{glosario/entradas_glosario}
% \addcontentsline{toc}{chapter}{Glosario}
% \printglossary

\end{document}
